% 摘要文件：只负责摘要页面的版式。
% 摘要的具体文字内容由 main.tex 里的宏提供，方便集中修改。

% 前置部分单独起页，并切换为罗马页码
\cleardoublepage
\pagenumbering{Roman}

% 中文摘要标题
\begin{center}
    {\zihao{-3}\songti\bfseries 中文摘要}
\end{center}

% 中文摘要主体框
\begin{tcolorbox}[
    enhanced,
    sharp corners,
    colback=white,
    colframe=black,
    width=\textwidth,
    height fill,
    boxrule=0.5pt,
    left=1.5em, right=1.5em,
    top=1em, bottom=1em,
    before skip=0pt,
    after skip=0pt
]
    % 摘要正文采用 1.5 倍行距
    \begin{spacing}{1.5}
        % 摘要标签固定，正文内容来自主文件定义的宏
        \noindent {\zihao{-4}\songti\bfseries 摘要：}
        {\zihao{-4}\songti \ChineseAbstractText}

        % 摘要与关键词之间保留空白
        \vspace{2\baselineskip}

        % 关键词行
        \noindent {\zihao{-4}\songti\bfseries 关键词：}
        {\zihao{-4}\songti \ChineseAbstractKeywords}
    \end{spacing}
\end{tcolorbox}

% 英文摘要单独分页
\clearpage

% 英文摘要标题
\begin{center}
    {\zihao{-3}\songti\bfseries English Abstract}
\end{center}

% 英文摘要主体框
\begin{tcolorbox}[
    enhanced,
    sharp corners,
    colback=white,
    colframe=black,
    width=\textwidth,
    height fill,
    boxrule=0.5pt,
    left=1.5em, right=1.5em,
    top=1em, bottom=1em,
    before skip=0pt,
    after skip=0pt
]
    \begin{spacing}{1.5}
        % 英文摘要正文
        \noindent {\zihao{-4}\bfseries Abstract:}
        {\zihao{-4}\songti \EnglishAbstractText}

        % 英文关键词
        \vspace{2\baselineskip}

        \noindent \textbf{Keywords:} \EnglishAbstractKeywords
    \end{spacing}
\end{tcolorbox}

% 摘要结束后换页，后面接目录
\clearpage